% ===== UNIT 25: Hyper-parameter selection and selection bias (5.7 part 1) =====
% Status: REWRITTEN  (change to REWRITTEN / REVIEWED when done; see CHECKLIST.md)
% Rules: style rewrite only. Keep every \cite, \ref, \label, number, equation and table value as is. Change the wording, not the content.
\subsection{Hyperparameter selection, held-out validation, and significance}
\label{sec:holdout}
We selected all post-processing hyperparameters by grid search on the VeRi-776 query/gallery split. These are AQE $k$, the $k$-reciprocal re-ranking triple $(k_1,k_2,\lambda)$ and the fusion weight $w_1$. This is a possible source of optimistic bias, and we state it openly. To be clear: every headline number in this paper uses these split-tuned values. The validation analyses below measure and bound the optimism of that choice. They are audits of the headline configuration, not its source.

Three observations limit that bias. First, the selection surface is flat, not peaked. The fused mAP stays above $93.0\%$ across the entire range $w_1\in[0.2,0.5]$ (Figure~\ref{fig:weight}), so the reported result does not sit on a sharp test-set optimum. Second, the re-ranking and AQE settings are close to established defaults. The values $k_1{=}80,k_2{=}15,\lambda{=}0.2$ lie within the ranges recommended by Zhong et al.~\cite{zhong2017rerank}, and AQE $k{=}3$ is a standard small neighbourhood. So they are not specially fitted to VeRi-776. Third, the test set (200 identities, 1\,678 queries) is large compared with the small number of scalar hyperparameters we selected. This limits how much we can overfit.
